% ================================================================
% FULL SYSTEMS AGENDA (saved for future use as standalone vision paper)
% Originally Appendix in "Determination Provenance: From Ambiguity to Algebra"
% ================================================================

\section{Systems Agenda}
\label{app:systems-lineage}
\label{app:certificates}

The theory developed in this paper appears to unlock practical
relevance for data provenance in settings---transactional databases,
distributed services, concurrent pipelines---where formal provenance
has had less impact on mainstream tracing and debugging infrastructure
than its theoretical power might suggest.
This appendix sketches a research agenda for systems work building on
determination provenance, mapping the framework's three levels of
structure to concrete capabilities beyond what current lineage
infrastructure provides.

Systems lineage---as practiced in distributed
tracing~\cite{sigelman2010dapper}, workflow
provenance~\cite{seltzer2005provenance}, and pipeline
auditing~\cite{bates2015trustworthy}---records \emph{which} services,
operations, or data sources contributed to an observed result.
This is lineage in the sense of Cui et al.~\cite{cui2000lineage}:
set-valued input attribution without algebraic structure.
It answers ``which inputs contributed?'' but not
``how did they combine?'' (no semiring), ``would the result
survive a different execution?'' (no counterfactuals), or ``which
assumptions were necessary?'' (no determination).
The result is verbose (full execution traces are retained) yet
expressively limited (only positive, single-execution queries are
supported).

Determination provenance addresses both deficiencies.
We map the framework's three levels of structure
(within a determination, across determinations, across specifications)
to concrete capabilities, then develop two orthogonal
dimensions of parsimony that make the theory practical.

\paragraph{Observation functions beyond relational databases.}
The body defines an observation function
$\mathsf{obs}: O \to \mathrm{Inst}(\mathbf{S})$ mapping outcomes to
relational instances.
In systems settings, observables are often partially-ordered,
semi-structured records (distributed traces, replicated logs, event
streams).
The framework applies whenever provenance queries over such records
can be expressed in a language equivalent to positive RA --- for
example, the positive fragment of a semi-structured query language
(select, filter, and join over JSON log entries, without negation or
aggregation) is RA-equivalent and natural for log analysis.

\subsection{Level 1: Repeatability from a Single Determination}

A resolving determination records the non-deterministic commitments
that resolved the specification to a function.
Once a determination is fixed, the remaining execution is deterministic:
given the history and the commitments, the resolved outcome can be
reconstructed by re-execution.
(If entailment steps involve non-determinism---e.g., hash-based
tie-breaking---the determination must be augmented with random seeds
for those steps; we assume deterministic entailment throughout.)

This is already a compression over full-trace lineage: rather than
recording every intermediate state, the system retains only the
commitment sequence.
All derivational detail can be reconstructed on demand.
Once reconstructed, classical semiring provenance applies, answering
positive provenance queries (why-provenance, how-provenance) with full
algebraic structure.

More broadly, determination provenance serves as a \emph{query layer}
on top of existing logs.
Full traces are retained as today; the determination (the
non-deterministic commitments and their structure) is extracted as a
\emph{semantic index} over the log.
Diagnostic queries---robustness, counterfactuals, root-cause
identification---are evaluated against the determination structure
($O(k)$ commitments) rather than by scanning raw events ($O(ne)$
log records).
The algebraic structure enables queries that log-scanning cannot answer
at all (e.g., ``under what fraction of admissible schedules does the
SLA hold?''), and answers those it can more efficiently by routing
through the determination semiring rather than correlating events.

\begin{example}[Distributed tracing {\`a} la Dapper~\cite{sigelman2010dapper}]
  \label{ex:dapper}
  A web-search request spawns a tree of spans: the frontend ($F$)
  calls a cache ($C$) and an ad server ($A$) in parallel; the cache
  misses and calls a database ($D$).
  Each span has client-send/server-receive events on different hosts
  with unsynchronized clocks.
  Causal constraints ($\mathsf{cs}_X \prec \mathsf{sr}_X$,
  $\mathsf{ss}_X \prec \mathsf{cr}_X$) are known, but the relative
  wall-clock durations of $C$ and $A$ are not---multiple orderings are
  admissible.

  A resolving determination commits to clock alignments
  ($\varphi_{\mathrm{align}(h_i,h_j)}$) between hosts, fixing a single
  execution timeline.
  Once fixed, critical-path analysis is deterministic: either $C$
  (cache miss + DB) or $A$ (ad auction) is on the critical path.

  Determination provenance enables queries that Dapper cannot answer:
  \begin{itemize}[nosep]
    \item \emph{Is span $D$ robustly on the critical path?}
          If $\mathrm{supp}(P(\mathrm{cp}=D)) = \mathcal{D}$, the
          database is the bottleneck under every admissible clock
          alignment.
    \item \emph{Under which clock assumptions does the ad server
          dominate latency?}
          The support $\mathrm{supp}(P(\mathrm{cp}=A))$ identifies
          exactly the determinations where $A$ is critical.
    \item \emph{Would migrating the cache to a closer datacenter
          change the critical path?}
          A cross-specification question: compare supports under the
          current spec vs.\ a spec with tighter latency bounds on $C$.
  \end{itemize}
\end{example}

\subsection{Level 2: Robustness, Tail Bounds, and Scheduler Design}

Supports across the space of resolving determinations enable queries
that systems lineage cannot express.

\paragraph{Robustness and fragility.}
The ratio $|\mathrm{supp}(P(t))| / |\mathcal{D}|$ measures how
\emph{fragile} an outcome is: a result with support ratio $0.99$ is
nearly robust (almost any schedule produces it); one with ratio $0.01$
depends on a very specific interleaving.
Boolean robustness ($\mathrm{supp} = \mathcal{D}$) is the extreme case.

\paragraph{Structural tail bounds.}
Consider a database with $n$ concurrent transactions and a latency
predicate $\ell(T_Q) < \tau$ (``query $T_Q$ completes within $\tau$
ms'').
The support of this predicate---the set of serializations under which
the SLA is met---gives a \emph{structural} tail-bound estimate:
$|\mathrm{supp}(\ell < \tau)| / |\mathcal{D}|$ is the fraction of
admissible schedules satisfying the SLA, without empirical sampling.
If this fraction is $0.95$, the system has a structural p95 guarantee.

\paragraph{Diagnosis via determination responsibility.}
When the tail is violated, determination responsibility
quantifies each commitment's contribution to the violation.
Rather than merely identifying which commitments are shared by all
violating determinations (a set-intersection test), responsibility
assigns each commitment a Shapley value measuring its marginal effect
on the SLA predicate.
A commitment with responsibility $0.73$ is the dominant root cause;
one with responsibility $0.02$ is incidental.
When the conflict structure has bounded treewidth---as in partitioned
workloads or tree-shaped service dependencies---these values are
efficiently computable.

\paragraph{Scheduler design.}
A scheduler that avoids high-responsibility commitments (e.g., never
serializes $T_i$ before $T_j$ on the contested object) eliminates the
tail without changing the isolation level or the workload.
This is a new optimization objective: not ``minimize average latency''
(a single-determination metric) but ``minimize the maximum
determination responsibility for the SLA-violation predicate''
(a distributional metric over the determination semiring).
Baccaert et al.~\cite{baccaert2026makespan} optimize makespan within a
single determination; the determination semiring enables optimization
\emph{across} determinations.

\subsection{Level 3: Cross-Specification Comparison}

Supports across variant specifications answer cross-specification
questions: does an outcome that holds under one specification also hold
under another?
Examples include isolation-level comparison (does a transaction's result
survive a move from SI to
serializability?~\cite{vandevoort2025mixed}), semantic-change queries
(would a $\mbox{Datalog}^\neg$ result survive a change from stable to well-founded
semantics?), and specification tightening (which results hold
when migrating to a stricter consistency model?).
These require comparing determination semirings of different
specs---a capability entirely absent from systems lineage.

\subsection{Parsimony via Truncation (Filtration)}
\label{sec:parsimony-truncation}

The filtration provides a structural answer to ``which layers do I need?''
If a query family $\mathcal{Q}$ depends only on tuples whose supports
lie in $\mathcal{F}_k$, then layers $k{+}1, \ldots, d$ are irrelevant:
the determination can be truncated at depth $k$ without affecting any
answer in $\mathcal{Q}$.

Query-relative depth makes this precise: a tuple with
$\mathrm{qdepth}(t) = k$ is fully explained by the first $k$ layers.
A system that retains only the first $k$ layers of each determination
can answer all queries whose inputs have depth $\le k$---and positive
relational algebra cannot increase depth beyond its inputs.

This is \emph{vertical} parsimony: discard layers that the query family
does not need.

\subsection{Parsimony via Compression (Certificates)}
\label{sec:parsimony-compression}

Within a layer that \emph{is} retained, the full commitment may carry
more detail than the query requires.
A \emph{certificate} is a sound over-approximation of a commitment's
effect: it may admit outcomes that the actual commitment excludes, but
excludes none that the commitment admits.

\begin{definition}[Certificate]
  A certificate for commitment $\varphi$ is a predicate
  $C : O \to \{\mathsf{true}, \mathsf{false}\}$ such that for all $H$,
  $\Spec(H \cdot \varphi) \subseteq \{o \in \Spec(H) \mid C(o)\}$.
\end{definition}

\begin{definition}[Sufficient certificate]
  A certificate $C$ is \emph{sufficient for a query family
  $\mathcal{Q}$} if for every $q \in \mathcal{Q}$ and every history $H$,
  $q(\Spec(H \cdot \varphi)) = q(\Spec(H \cdot C))$---that is,
  replacing $\varphi$ with $C$ yields the same answer.
\end{definition}

\begin{theorem}[Query monotonicity]
  If $\mathcal{Q}_1 \subseteq \mathcal{Q}_2$, every certificate
  sufficient for $\mathcal{Q}_2$ is sufficient for $\mathcal{Q}_1$.
\end{theorem}

\noindent
This is \emph{horizontal} parsimony: within each retained layer,
coarsen the representation to the minimum fidelity the query family
demands.

\begin{example}[Certificates in practice]
  Under serializability, a certificate for derivational provenance need
  only record the ordering predicates relevant to the queried tuple's
  conflict neighborhood.
  Answering a cross-specification query (``would this result survive
  under a different isolation level?'') requires a strictly stronger
  certificate that preserves the full conflict-graph structure.
  In distributed tracing, a certificate for critical-path analysis
  retains only the clock-alignment and ordering commitments on the
  path itself---not the full trace.
\end{example}

\paragraph{The two dimensions compose.}
For a given query family: first truncate (filtration identifies the
deepest relevant layer), then compress each retained layer (certificates
identify the minimum per-layer fidelity).
Together they define a well-parameterized design space for parsimonious
systems lineage.
Fixing the parameters---choosing the query family, computing minimal
certificates, implementing adaptive retention---is the systems work
that remains.
